% GENERATED FILE. Edit canonical lessons, then run npm run generate:books.
% canonical-source-hash: fnv1a64:b8fb97e541d307d4
% canonical-lessons: MR-C436-vatavaran, MR-C436-usnata, MR-C436-ukada, MR-C436-damat, MR-C436-talapne

\chapter{Heat and Humidity}
\label{ch:mr-heat-and-humidity}
\hlchaptermodality{436}

\begin{chapteropening}
\textbf{By the end of this chapter:} \emph{I can talk about the weather, great heat, sultry days and the blazing sun.}

Complete the last lesson of chapter 436: I can talk about the weather, great heat, sultry days and the blazing sun.
\end{chapteropening}

\section[\texorpdfstring{\mr{वातावरण} (vātāvaraṇ)}{वातावरण (vātāvaraṇ)}]{\mr{वातावरण} (vātāvaraṇ) — the atmosphere; the weather}

\label{lesson:MR-C436-vatavaran}



\begin{quote}
Before the new one: say the Marathi for to dry up, then the Marathi for to go bad.
\end{quote}

\subsection*{You'll want to know: \mr{वातावरण}}
\textbf{\mr{वातावरण}} — \emph{vātāvaraṇ} — \textquotedblleft{}the atmosphere; the weather\textquotedblright{}.

\begin{quote}
\mr{आज} \mr{वातावरण} \mr{ढगाळ} \mr{आहे}.
\end{quote}

\begin{quote}
\emph{āj vātāvaraṇ ḍhagāḷ āhe.}
\end{quote}

\begin{quote}
Today the weather is cloudy.
\end{quote}

Neuter.

Also the mood of a place: \mr{घरात} \mr{आनंदी} \mr{वातावरण} \mr{होतं}.

\begin{cousinweb}
\textbf{\mr{वातावरण}} (\emph{vātāvaraṇ}) is a modern word built from Sanskrit \emph{vāta}, \textquotedblleft{}wind, air\textquotedblright{}, and \emph{āvaraṇa}, \textquotedblleft{}a covering\textquotedblright{} (Monier-Williams): the cover of air.
\end{cousinweb}

\begin{grammarlens}[title={What you've built}]
One more word for getting around.
\end{grammarlens}

\subsection*{Your turn}
\begin{itemize}
\raggedright
  \item \emph{Say it:} \emph{vātāvaraṇ}
  \item \emph{Say it:} \emph{vātāvaraṇ}, once more
  \item \emph{Say it:} \emph{āj vātāvaraṇ ḍhagāḷ āhe.} — Today the weather is cloudy.
  \item \emph{Recall:} say the Marathi for to dry up, then the Marathi for to go bad, then say \emph{vātāvaraṇ} again
\end{itemize}

\begin{tcolorbox}[breakable,colback=teal!4,colframe=teal!35!black,title={Before you move on}]
What does \mr{वातावरण} mean? (\textquotedblleft{}The atmosphere; the weather\textquotedblright{}.) Say it once more.
\end{tcolorbox}
\section[\texorpdfstring{\mr{उष्णता} (uṣṇatā)}{उष्णता (uṣṇatā)}]{\mr{उष्णता} (uṣṇatā) — great heat; heat}

\label{lesson:MR-C436-usnata}



\begin{quote}
Before the new one: say the Marathi for to go bad, then the Marathi for the atmosphere.
\end{quote}

\subsection*{You'll want to know: \mr{उष्णता}}
\textbf{\mr{उष्णता}} — \emph{uṣṇatā} — \textquotedblleft{}great heat; heat\textquotedblright{}.

\begin{quote}
\mr{दुपारी} \mr{उष्णता} \mr{खूप} \mr{वाढते}.
\end{quote}

\begin{quote}
\emph{dupārī uṣṇatā khūp vāḍhte.}
\end{quote}

\begin{quote}
In the afternoon the heat rises a lot.
\end{quote}

Feminine.

A formal word, also heat in science.

\begin{cousinweb}
\textbf{\mr{उष्णता}} (\emph{uṣṇatā}) is a tatsama, taken from Sanskrit \emph{uṣṇatā}, \textquotedblleft{}heat\textquotedblright{}, from \emph{uṣṇa}, \textquotedblleft{}hot\textquotedblright{} (Monier-Williams).
\end{cousinweb}

\begin{grammarlens}[title={What you've built}]
One more word for getting around.
\end{grammarlens}

\subsection*{Your turn}
\begin{itemize}
\raggedright
  \item \emph{Say it:} \emph{uṣṇatā}
  \item \emph{Say it:} \emph{uṣṇatā}, once more
  \item \emph{Say it:} \emph{dupārī uṣṇatā khūp vāḍhte.} — In the afternoon the heat rises a lot.
  \item \emph{Recall:} say the Marathi for to go bad, then the Marathi for the atmosphere, then say \emph{uṣṇatā} again
\end{itemize}

\begin{tcolorbox}[breakable,colback=teal!4,colframe=teal!35!black,title={Before you move on}]
What does \mr{उष्णता} mean? (\textquotedblleft{}Great heat; heat\textquotedblright{}.) Say it once more.
\end{tcolorbox}
\section[\texorpdfstring{\mr{उकाडा} (ukāḍā)}{उकाडा (ukāḍā)}]{\mr{उकाडा} (ukāḍā) — sultry heat}

\label{lesson:MR-C436-ukada}



\begin{quote}
Before the new one: say the Marathi for the atmosphere, then the Marathi for great heat.
\end{quote}

\subsection*{You'll want to know: \mr{उकाडा}}
\textbf{\mr{उकाडा}} — \emph{ukāḍā} — \textquotedblleft{}sultry heat\textquotedblright{}.

\begin{quote}
\mr{पावसाआधी} \mr{खूप} \mr{उकाडा} \mr{होतो}.
\end{quote}

\begin{quote}
\emph{pāvsāādhī khūp ukāḍā hoto.}
\end{quote}

\begin{quote}
Before the rain it gets very sultry.
\end{quote}

Masculine.

Heavy, sticky heat; \mr{उकाडा} \mr{होणे}.

\begin{cousinweb}
\textbf{\mr{उकाडा}} (\emph{ukāḍā}) is inherited, from the family of Sanskrit \emph{utkvathati}, \textquotedblleft{}boils out\textquotedblright{}, where CDIAL lists Marathi \emph{ukaḍṇẽ}, \textquotedblleft{}to be hot (of weather)\textquotedblright{} (CDIAL 1741).
\end{cousinweb}

\begin{grammarlens}[title={What you've built}]
One more word for getting around.
\end{grammarlens}

\subsection*{Your turn}
\begin{itemize}
\raggedright
  \item \emph{Say it:} \emph{ukāḍā}
  \item \emph{Say it:} \emph{ukāḍā}, once more
  \item \emph{Say it:} \emph{pāvsāādhī khūp ukāḍā hoto.} — Before the rain it gets very sultry.
  \item \emph{Recall:} say the Marathi for the atmosphere, then the Marathi for great heat, then say \emph{ukāḍā} again
\end{itemize}

\begin{tcolorbox}[breakable,colback=teal!4,colframe=teal!35!black,title={Before you move on}]
What does \mr{उकाडा} mean? (\textquotedblleft{}Sultry heat\textquotedblright{}.) Say it once more.
\end{tcolorbox}
\section[\texorpdfstring{\mr{दमट} (damaṭ)}{दमट (damaṭ)}]{\mr{दमट} (damaṭ) — humid}

\label{lesson:MR-C436-damat}



\begin{quote}
Before the new one: say the Marathi for great heat, then the Marathi for sultry heat.
\end{quote}

\subsection*{You'll want to know: \mr{दमट}}
\textbf{\mr{दमट}} — \emph{damaṭ} — \textquotedblleft{}humid\textquotedblright{}.

\begin{quote}
\mr{समुद्राजवळची} \mr{हवा} \mr{दमट} \mr{असते}.
\end{quote}

\begin{quote}
\emph{samudrājavaḷcī havā damaṭ aste.}
\end{quote}

\begin{quote}
The air near the sea is humid.
\end{quote}

It does not change: \mr{दमट} \mr{हवा}, \mr{दमट} \mr{कपडे}.

\begin{cousinweb}
\textbf{\mr{दमट}} (\emph{damaṭ}) is a Marathi adjective for damp, sticky air; its origin is not settled.
\end{cousinweb}

\begin{grammarlens}[title={What you've built}]
One more word for getting around.
\end{grammarlens}

\subsection*{Your turn}
\begin{itemize}
\raggedright
  \item \emph{Say it:} \emph{damaṭ}
  \item \emph{Say it:} \emph{damaṭ}, once more
  \item \emph{Say it:} \emph{samudrājavaḷcī havā damaṭ aste.} — The air near the sea is humid.
  \item \emph{Recall:} say the Marathi for great heat, then the Marathi for sultry heat, then say \emph{damaṭ} again
\end{itemize}

\begin{tcolorbox}[breakable,colback=teal!4,colframe=teal!35!black,title={Before you move on}]
What does \mr{दमट} mean? (\textquotedblleft{}Humid\textquotedblright{}.) Say it once more.
\end{tcolorbox}
\section[\texorpdfstring{\mr{तळपणे} (taḷapṇe)}{तळपणे (taḷapṇe)}]{\mr{तळपणे} (taḷapṇe) — to blaze}

\label{lesson:MR-C436-talapne}



\begin{quote}
Before the new one: say the Marathi for sultry heat, then the Marathi for humid.
\end{quote}

\subsection*{You'll want to know: \mr{तळपणे}}
\textbf{\mr{तळपणे}} — \emph{taḷapṇe} — \textquotedblleft{}to blaze\textquotedblright{}.

\begin{quote}
\mr{दुपारी} \mr{सूर्य} \mr{तळपत} \mr{होता}.
\end{quote}

\begin{quote}
\emph{dupārī sūrya taḷpat hotā.}
\end{quote}

\begin{quote}
At noon the sun was blazing.
\end{quote}

Past (he …): \mr{तो} \mr{तळपला}, \emph{to taḷaplā}.

For the sun and for a bright sword.

\begin{cousinweb}
\textbf{\mr{तळपणे}} (\emph{taḷapṇe}) is inherited: CDIAL lists Marathi \emph{taḷapṇẽ}, \textquotedblleft{}to be brandished\textquotedblright{}, under \emph{*taḍapphaḍ}, \textquotedblleft{}agitate\textquotedblright{} (CDIAL 5634).
\end{cousinweb}

\begin{grammarlens}[title={What you've built}]
One more word for getting around.
\end{grammarlens}

\subsection*{Your turn}
\begin{itemize}
\raggedright
  \item \emph{Say it:} \emph{taḷapṇe}
  \item \emph{Say it:} \emph{taḷapṇe}, once more
  \item \emph{Say it:} \emph{dupārī sūrya taḷpat hotā.} — At noon the sun was blazing.
  \item \emph{Recall:} say the Marathi for sultry heat, then the Marathi for humid, then say \emph{taḷapṇe} again
\end{itemize}

\begin{tcolorbox}[breakable,colback=teal!4,colframe=teal!35!black,title={Before you move on}]
What does \mr{तळपणे} mean? (\textquotedblleft{}To blaze\textquotedblright{}.) Say it once more.
\end{tcolorbox}
